\cotefile{157-2}
\scanpage{398}\mspage{137}
12. Retour sur les ``catégories de chemins'' \deleted{\unsure{d}}dans $\mathrm{Cat}$, comme
catégorie de modèles pour $\mathrm{Hot}$.

Je \unsure{vois|vais} \unsure{actuellement} deux applications intéressantes
du formalisme des chemins dans les catégories.
L'une, qui m'a surtout forcé la main, est
pour pouvoir construire des $\mathrm{Hom}\,\mathrm{cat}$, \uline{$\mathrm{Hom}$
externes} à valeurs dans $\mathrm{Hot}$, pour tout \inserted{pré}dérivateur
$\mathbb D$ défini par un couple $(M,\Sigma)$ (et je
ne connais pas \unsure{de|des} prédérivateur qui ne
soit défini ainsi.

L'autre application est pour avoir dans $(\mathrm{Cat}$,
une \uline{construction} \uline{fonctorielle}, en termes
d'un objet $X$ quelconque ou d'une \deleted{\unsure{$f\colon Y$|$f\colon X$}}
flèche $X\unsure{\text{---}|\longrightarrow}Y$, de l'analogue du formalisme
de Cartan-Serre sur les espaces de
chemins. Dans ce cadre, il est entendu
qu'on prend toujours $\Sigma_X=\mathrm{Fl}(X)$, ensemble de
\uline{toutes} les flèches de $X$. Pour construire
fonctoriellement en $X$ de $(\mathrm{Cat})$
\diagram{\begin{tikzcd}[ampersand replacement=\&]
\underline{\mathrm{Ch}}_{\infty}(X) \arrow[d] \\
X\times X
\end{tikzcd}}
il faut avoir fixé un syst. projectif
\[
\underline t=(t_n)_{\unsure{n\in\mathbb N|\in\mathbb N}}
\]
% archive sheet number: 397
